\documentclass[../main.tex]{subfiles}
\begin{document}

\chapter{Many Cores}
\lessoninfo{
  video={Lesson 22, videos 1--19},
  textbook={H\&P \cite{hennessy2017} §1.5, §5.4, App.~F},
  keywords={many-core processor, coherence traffic, network on chip, mesh, torus, tile, off-chip traffic, distributed last-level cache, on-chip directory, partial directory, power budget, turbo mode, logical processor},
}

Lessons 18--21 developed multiprocessors whose cores typically share a bus
and a last-level cache, keep their caches coherent by snooping on the bus
(or, in a machine without a common bus, through a directory), coordinate
with locks, and rely on a consistency model for shared data.  Every
technology generation fits more cores on a chip, however, and mechanisms
that serve four or eight cores well break down when a chip has dozens or
hundreds of them.  This lesson examines the challenges
that such chips face---coherence traffic, off-chip traffic, the size of the
coherence directory, the shared power budget and the placement of threads
by the operating system---and the solution the lecture gives for each.

\section{Coherence traffic}
\label{sec:l22-coherence}

A \term{many-core processor}[メニーコアプロセッサ] is a multicore processor
with many cores on one chip: dozens to hundreds rather than two to
eight.\source{Lesson 22, videos 1--2; H\&P ch.~5.}  The first challenge of
such a chip concerns cache coherence (Lesson 19).  Every write to a shared
location that other cores have cached causes \emph{coherence traffic}:
invalidations sent to the other caches, and coherence misses when those
cores access the block again.  With more cores the chip performs more writes
per second in total, so the rate of coherence messages grows with the
number of cores.  On a shared bus this traffic soon becomes the bottleneck:
\begin{itemize}
  \item a bus carries one transfer at a time, so the cores wait for each
    other's coherence messages and cache misses;
  \item a bus that connects more cores is longer and has more devices
    attached, so each transfer is slower as well.
\end{itemize}

A many-core chip therefore needs two things: an on-chip interconnect whose
capacity grows with the number of cores, and a coherence mechanism that
does not rely on every cache observing every operation on a single bus.
The second is directory-based coherence (\cref{sec:l19-directory}), which
delivers each message only to the caches that hold the block, instead of
broadcasting it to all cores.

\section{Network on chip}
\label{sec:l22-noc}

The shared bus is replaced by a network.\source{Lesson 22, videos 3--4;
H\&P App.~F; \cite{dally2001}.}

\begin{definition}[Network on chip]
  A \term{network on chip}[ネットワークオンチップ] connects the cores of a
  chip by point-to-point links between routers, one router per core.  A
  message travels from router to router until it reaches its destination,
  and messages on different links travel at the same time.
\end{definition}

In the \term{mesh}[メッシュ] topology the routers form a two-dimensional
grid, and each router is linked to its neighbours to the north, south, east
and west (\cref{fig:l22-topologies}b).  Any core reaches any other core over
a path of several links, or \emph{hops}, and transfers that use different
links proceed at the same time, so the network as a whole carries many
times the traffic of a single link.  Two further properties suit the mesh to
a chip.
\begin{description}
  \item[Scaling] Adding cores enlarges the grid, and every added core adds
    links, so the capacity of the network grows with the number of cores,
    unlike that of a bus.  Every link connects two neighbouring cores and
    stays short however large the chip, whereas a bus must span all cores.
  \item[Planar layout] The links connect only neighbours and do not cross,
    so the mesh is easily laid out on the flat surface of a chip.
\end{description}
Such a chip is built by replicating a \term{tile}[タイル]: a core with its
private caches and its router, to which later sections add a part of the
structures that all cores share.

\begin{figure}[htbp]
  \centering
  % Interconnect topologies for many cores: (a) a shared bus, (b) a 4x4 mesh
% carrying two transfers at the same time on different links, (c) a 4x4
% torus, i.e. a mesh with wraparound links in every row and column.
% Links are drawn between node centres on the background layer, so the
% filled nodes hide the parts of the lines inside them.
\begin{tikzpicture}[x=1cm,y=1cm, font=\footnotesize\sffamily,
  nd/.style={draw, fill=ln-accent!15, minimum size=0.36cm, inner sep=0pt},
  lk/.style={ln-muted, line width=0.5pt},
  wrap/.style={ln-accent, line width=0.5pt, rounded corners=2pt},
  route/.style={line width=1.6pt, line cap=round, line join=round},
  paneltitle/.style={font=\footnotesize\sffamily, anchor=north}]
  % ---------------------------------------------------------------- (a) bus
  \begin{scope}
    \foreach \i in {0,...,3} {
      \node[nd] at (0.3+\i*0.8,2.3) {};
      \node[nd] at (0.3+\i*0.8,0.4) {};
    }
    \begin{scope}[on background layer]
      \foreach \i in {0,...,3}
        \draw[lk] (0.3+\i*0.8,2.3) -- (0.3+\i*0.8,0.4);
    \end{scope}
    \draw[line width=1.4pt] (-0.1,1.35) -- (3.1,1.35);
    \node[paneltitle] at (1.5,-0.6) {\iflangja{(a) バス}{(a) bus}};
  \end{scope}
  % --------------------------------------------------------------- (b) mesh
  \begin{scope}[xshift=4.0cm]
    \foreach \x in {0,...,3} \foreach \y in {0,...,3}
      \node[nd] at (\x*0.9,\y*0.9) {};
    \begin{scope}[on background layer]
      \foreach \k in {0,...,3} {
        \draw[lk] (\k*0.9,0) -- (\k*0.9,2.7);
        \draw[lk] (0,\k*0.9) -- (2.7,\k*0.9);
      }
      % two transfers in progress on disjoint links
      \draw[route, ln-caution] (0,2.7) -- (1.8,2.7) -- (1.8,1.8);
      \draw[route, ln-tip] (2.7,0) -- (2.7,0.9) -- (0.9,0.9);
    \end{scope}
    \node[paneltitle] at (1.35,-0.6) {\iflangja{(b) メッシュ}{(b) mesh}};
  \end{scope}
  % -------------------------------------------------------------- (c) torus
  \begin{scope}[xshift=8.1cm]
    \foreach \x in {0,...,3} \foreach \y in {0,...,3}
      \node[nd] at (\x*0.9,\y*0.9) {};
    \begin{scope}[on background layer]
      \foreach \k in {0,...,3} {
        \draw[lk] (\k*0.9,0) -- (\k*0.9,2.7);
        \draw[lk] (0,\k*0.9) -- (2.7,\k*0.9);
        % wraparound links: below each row, right of each column
        \draw[wrap] (0,\k*0.9) -- (-0.32,\k*0.9) -- (-0.32,\k*0.9-0.3)
          -- (3.02,\k*0.9-0.3) -- (3.02,\k*0.9) -- (2.7,\k*0.9);
        \draw[wrap] (\k*0.9,2.7) -- (\k*0.9,3.02) -- (\k*0.9+0.3,3.02)
          -- (\k*0.9+0.3,-0.32) -- (\k*0.9,-0.32) -- (\k*0.9,0);
      }
    \end{scope}
    \node[paneltitle] at (1.35,-0.6) {\iflangja{(c) トーラス}{(c) torus}};
  \end{scope}
\end{tikzpicture}

  \caption{Interconnects for many cores.  A bus carries one transfer at a
    time; the mesh carries the two highlighted transfers at the same time on
    different links; the torus adds a wraparound link between the first and
    the last router of every row and every column---only the link of the
    left column and the one of the bottom row are drawn.}
  \label{fig:l22-topologies}
\end{figure}

Other topologies shorten the paths at the cost of a more difficult layout.
A \term{torus}[トーラス] is a mesh with additional wraparound links between
the first and the last router of every row and every column
(\cref{fig:l22-topologies}c), which roughly halves the longest path; the
wraparound links, however, are long or cross other links.  The
\emph{flattened butterfly} and the \emph{hypercube} link each router to more
routers---to all routers of its own row and column in a flattened butterfly,
to $\log_{2}N$ routers in a hypercube---which shortens the paths further but
requires long, crossing links that are hard to build on a chip.

If the traffic were spread evenly over the links, a network of $L$ links,
each carrying up to $C$ messages per second, would carry
\begin{equation}
  \text{throughput} \;\approx\; \frac{L\cdot C}{h}
  \label{eq:l22-throughput}
\end{equation}
messages per second, where $h$ is the average number of hops per message,
because a message occupies a link on every hop.  A $k\times k$ mesh has
$L=2k(k-1)$ links, and for a destination drawn at random among the other
tiles the average distance is $h=2k/3$ hops, which grows only with the
square root of the number of cores $N=k^{2}$.

\Cref{eq:l22-throughput} is an average-load estimate that no real routing
reaches.  The load on a link is the sum of the rates of all messages routed
over it, and the network carries the offered traffic only while every link
stays below $C$, so the busiest link, not the average, sets the
limit.\caution{A mesh raises throughput but does not shorten latency: a
message to a distant tile passes many routers, each adding delay.}  When
messages are addressed all over the chip and the routers send each message
along one dimension and then the other, the links that cross the middle of
the mesh are the busiest, and they saturate at about $2kC$ messages per
second.  The capacity of a mesh therefore grows with $\sqrt{N}$ for such
traffic, and in proportion to $N$ only while messages stay among
neighbouring tiles; a bus, by contrast, never carries more than one link's
worth of traffic however many cores it connects.  Since the capacity per
core falls as $1/\sqrt{N}$ once traffic crosses the chip, keeping a core's
data in its own tile matters (\cref{sec:l22-llc}).

\begin{example}
  \label{ex:l22-throughput}
  Let each link, and the bus, carry $C$ messages per second.  With 16 cores
  the bus carries $C$; a $4\times4$ mesh has 24 links and an average
  distance of $8/3\approx 2.67$ hops, for which \cref{eq:l22-throughput}
  estimates $24C/2.67\approx 9C$, while its middle links limit it to about
  $8C$.  With 64 cores the bus still carries at most $C$ (in practice less,
  since it is longer), while an $8\times8$ mesh has 112 links and an average
  distance of $16/3\approx 5.33$ hops, which the estimate turns into $21C$
  and the middle links into about $16C$: four times the cores, twice the
  capacity.  If all traffic went to neighbouring tiles ($h=1$), the two
  meshes would carry $24C$ and $112C$.
\end{example}

\section{Off-chip traffic and the distributed LLC}
\label{sec:l22-llc}

The second challenge is the traffic that leaves the chip.\source{Lesson 22,
videos 5--7.}  Each core has its own caches, and the miss rate per core does
not change with the number of cores, so a chip with more cores sends
proportionally more requests to main memory.  These requests pass through
the pins of the chip; the number of pins also grows from generation to
generation, but far more slowly than the number of cores, so the memory
interface becomes a bottleneck.  The number of memory requests \emph{per
core} must therefore go down: more accesses must be served on the chip, by
a last-level cache that all cores share and whose size grows roughly in
proportion to the number of cores.  A single large LLC, however, is slow,
since a larger cache has a longer hit time (Lesson 12), and it has a single
entry point through which the L2 misses of all cores must pass, which is a
bottleneck of its own.

\begin{definition}[Distributed last-level cache]
  A \term{distributed last-level cache}[分散最終レベルキャッシュ] is
  logically a single cache shared by all cores---each block is stored in it
  at most once, not once per core---but is physically divided into
  \emph{slices}, one slice in each tile (\cref{fig:l22-llc}).
\end{definition}

The size of the LLC is the number of cores times the size of a slice, so it
grows automatically with the number of cores, while each slice remains
small and fast and the requests of the cores are spread over many slices.

\begin{figure}[htbp]
  \centering
  % A many-core chip of eight tiles connected by a mesh.  Each tile holds a
% core, its private L1 and L2 caches, one slice of the distributed LLC, the
% part of the on-chip directory kept in that tile, and a router (R).  An L2
% miss of core 1 on a block of a set whose number ends in 110 goes to LLC
% slice 6.
\begin{tikzpicture}[x=1cm,y=1cm, font=\scriptsize\sffamily, >={Stealth[length=1.8mm]},
  tile/.style={draw=ln-rule, rounded corners=2pt},
  core/.style={draw, fill=ln-accent!15, minimum width=1.5cm, minimum height=0.34cm, inner sep=0pt},
  box/.style={draw, fill=ln-box, minimum height=0.34cm, inner sep=0pt},
  rt/.style={draw, circle, fill=white, minimum size=0.36cm, inner sep=0pt, font=\tiny\sffamily},
  lk/.style={ln-muted, line width=0.6pt},
  hot/.style={ln-caution, line width=1.5pt}]
  \foreach \i in {0,...,7} {
    \pgfmathtruncatemacro{\col}{mod(\i,4)}
    \pgfmathtruncatemacro{\row}{div(\i,4)}
    \begin{scope}[shift={(\col*2.65,-\row*2.6)}]
      \draw[tile] (-1.1,-1.35) rectangle (1.2,0.95);
      \node[core] (core\i) at (-0.2,0.62) {\iflangja{コア}{core}~\i};
      \node[box, minimum width=0.72cm] (l1\i) at (-0.59,0.18) {L1};
      \node[box, minimum width=0.72cm] (l2\i) at (0.19,0.18) {L2};
      \node[box, minimum width=1.5cm] (llc\i) at (-0.2,-0.26) {LLC~\i};
      \node[box, minimum width=1.5cm] (dir\i) at (-0.2,-0.70) {\iflangja{ディレクトリ}{directory}~\i};
      \coordinate (rp\i) at (0.95,-1.12);
    \end{scope}
  }
  % mesh links between routers (drawn before the routers)
  \foreach \i/\j in {0/1,1/2,2/3,4/5,5/6,6/7,0/4,1/5,2/6,3/7}
    \draw[lk] (rp\i) -- (rp\j);
  % route of the request: R1 -> R2 -> R6
  \draw[hot] (rp1) -- (rp2) -- (rp6);
  \fill[ln-caution!20] (llc6.south west) rectangle (llc6.north east);
  \draw (llc6.south west) rectangle (llc6.north east);
  \node at (llc6) {LLC~6};
  \foreach \i in {0,...,7} \node[rt] (r\i) at (rp\i) {R};
  \draw[->, hot, line width=0.9pt] (l21.east) -| (r1.north);
  \draw[->, hot, line width=0.9pt] (r6.west) -| ([xshift=0.2cm]llc6.east) -- (llc6.east);
\end{tikzpicture}

  \caption{Eight tiles connected by a mesh.  Each tile holds a core, its L1
    and L2 caches, one slice of the distributed LLC, a part of the on-chip
    directory (\cref{sec:l22-directory}) and a router (R).  An L2 miss of
    core~1 on a block whose set number ends in binary 110 is sent to LLC
    slice~6.}
  \label{fig:l22-llc}
\end{figure}

On an L2 miss the core must send its request to the one slice that may hold
the block, and it must be able to compute that slice from the address alone.
The lecture distributes the blocks among the tiles in round-robin fashion,
in one of two ways.
\begin{description}
  \item[By cache index] The lowest bits of the set number select the tile, so
    consecutive sets go to consecutive tiles.  The blocks are spread evenly,
    but without regard to which core uses them: the blocks of a thread are
    scattered over all tiles, and most of its LLC accesses travel to a
    distant tile.
  \item[By page number] Bits of the physical page number select the tile, so
    all blocks of a page are in the same slice.  The operating system
    decides which physical frame each page receives; it can give a thread
    frames whose slices lie in or near the thread's own tile, and most LLC
    accesses stay local---the same idea as keeping a thread's memory on its
    own node in a NUMA machine (\cref{sec:l18-distributed}).
\end{description}

\begin{example}
  \label{ex:l22-llc}
  A chip has 32 tiles with a \SI{1}{MB} LLC slice each, so its LLC holds
  \SI{32}{MB}.  With \SI{64}{B} blocks and 8-way set associativity, the LLC
  has $2^{25}/(2^{6}\cdot 2^{3})=2^{16}$ sets, of which each slice holds
  $2^{11}=2048$.  If the tile is selected by cache index, the 16-bit set
  number consists of address bits 6--21: its lowest 5 bits (address bits
  6--10) select one of the 32 tiles, the remaining 11 bits (address bits
  11--21) select the set within that slice, and the tag consists of the
  address bits from bit 22 up, as in any cache with $2^{16}$ sets.
\end{example}

\section{The on-chip directory}
\label{sec:l22-directory}

Directory-based coherence brings a challenge of its own.\source{Lesson 22,
videos 8--11; H\&P §5.4; \cite{agarwal1988}.}  The directory of
\cref{sec:l19-directory} has an entry for every memory block, and each entry
has one presence bit per core plus an exclusive bit.  Main memory of many
gigabytes consists of hundreds of millions of blocks, and a directory with
an entry for each of them does not fit on a chip.

\begin{example}
  \label{ex:l22-directory-size}
  A 64-core system has \SI{32}{GB} of memory divided into \SI{64}{B} blocks,
  that is, $2^{35}/2^{6}=2^{29}\approx 537$ million blocks.  At 65 bits per
  entry the full directory needs about $3.5\times10^{10}$ bits, or
  \SI{4.1}{GB}---about \SI{13}{\percent} of main memory itself.
\end{example}

Two decisions adapt the directory to the chip.  The first is where an entry
is kept.  The directory is divided among the tiles, and the entry of a block
resides in the tile that holds the block's LLC slice, its \emph{home} tile.
A core whose private caches miss must contact that tile anyway to search the
LLC, so it reaches the directory entry without an additional message; no
tile becomes a central bottleneck either.  The second decision, the one that
makes the directory small enough, is which blocks have an entry at all.

\begin{definition}[On-chip directory]
  An \term{on-chip directory}[オンチップディレクトリ] is a coherence
  directory kept on the chip: it is distributed over the tiles, the entry of
  each block residing in the tile of the block's LLC slice, and it is a
  \term{partial directory}[部分ディレクトリ], which has a fixed number of
  entries in each tile and holds entries only for some of the blocks.
\end{definition}

An entry is needed only for a block that is held by at least one private
cache (L1 or L2), that is, whose entry would have at least one presence bit
set.  A block that is only in the LLC or in main memory would have an entry
of zeros, and a missing entry says the same thing.  The directory therefore
allocates an entry when a private cache obtains a block that has none, and
the entry is no longer needed once no private cache holds the block.  When a
new entry is needed but all entries of the tile are in use, the directory
replaces one:
\begin{enumerate}
  \item it selects a victim entry, for example the least recently used one;
  \item it sends an invalidation to every tile whose presence bit is set in
    the victim entry; these caches discard their copies of the block (a
    modified copy is first written back);
  \item it uses the entry for the new block.
\end{enumerate}
The invalidated caches miss the next time they access the block, although
no other core has written to it.  Directory replacement thus causes a new
kind of cache miss, in addition to the misses of the 3C model (Lesson 14)
and the coherence misses of Lesson 19.  The more entries the directory
has, the rarer such misses become.\margin{The partial directory behaves like a
cache of directory entries, with its own replacement policy.}

\begin{example}
  \label{ex:l22-directory-trace}
  \Cref{tab:l22-directory-trace} traces the directory part of one tile with
  room for two entries in a chip with four cores, C0--C3.  Blocks A, B and C
  all have this tile as their home tile, and initially no private cache holds
  them.
\end{example}

\begin{table}[htbp]
  \centering\small
  \caption{A partial directory with two entries.  Each entry shows the
    block, the exclusive bit and the presence bits of C0--C3.}
  \label{tab:l22-directory-trace}
  \begin{tabular}{@{}clccl@{}}
    \toprule
    Step & Event & Entry 1 & Entry 2 & Directory action \\
    \midrule
    1 & C0 reads A  & \texttt{A 1 1000} & ---               & allocate entry 1 \\
    2 & C2 reads A  & \texttt{A 0 1010} & ---               & forward read to C0 \\
    3 & C1 writes B & \texttt{A 0 1010} & \texttt{B 1 0100} & allocate entry 2 \\
    4 & C3 reads C  & \texttt{C 1 0001} & \texttt{B 1 0100} & replace A: invalidate C0, C2 \\
    5 & C0 reads A  & \texttt{C 1 0001} & \texttt{A 1 1000} & replace B: invalidate C1 \\
    \bottomrule
  \end{tabular}
\end{table}

In step~4 the directory is full, and A, last used in step~2, is the least
recently used entry; C0 and C2 lose their copies of A.  The read of A by C0
in step~5 is a miss caused only by that replacement, and it in turn evicts B
from the directory and from C1, whose modified copy is written back to the
LLC.

\begin{keypoint}
  The distributed LLC and the on-chip directory both divide a shared
  structure among the tiles, so that it grows with the number of cores
  without a central bottleneck.  A block's LLC slice and its directory entry
  reside in the same tile, and the directory keeps entries only for blocks
  in private caches, at the price of occasional misses caused by directory
  replacement.
\end{keypoint}

\section{Sharing the power budget}
\label{sec:l22-power}

The third challenge is power.\source{Lesson 22, videos 12--14; H\&P §1.5.}
A chip can dissipate only as much power as its package and cooling remove,
and this power budget is shared by all its cores.  The more cores a chip
has, the less power each core receives, and the lower its supply voltage
and clock frequency must be.  The dynamic power of a core is proportional to
$V^{2}f$ (\cref{eq:active-power}), and the frequency a core reaches is
roughly proportional to its voltage, so the power grows roughly with $f^{3}$
(\cref{ex:l18-power}).  If $N$ cores share the budget of a single core, each
receives $1/N$ of it and runs at
\begin{equation}
  f_{N} \;=\; \frac{f_{1}}{\sqrt[3]{N}} ,
  \label{eq:l22-freq}
\end{equation}
where $f_{1}$ is the frequency of the single core.  This relation assumes
that the voltage can be lowered in proportion to the frequency; near its
lower limit, power falls only linearly with $f$, and dividing the budget
costs even more frequency.

\begin{marginfigure}
  \resizebox{\linewidth}{!}{% Performance of N cores sharing the power budget of one core, relative to
% the single core (P proportional to f^3, so each core runs at N^(-1/3) of
% the single-core frequency).  Curves: a program that is fully parallel,
% one with parallel fraction 0.9, and a single-threaded program.
\begin{tikzpicture}[x=0.75cm,y=0.8cm, font=\scriptsize\sffamily]
  \draw[ln-rule] (0,0) grid[xstep=1,ystep=1] (6,4);
  \draw[->] (0,0) -- (6.3,0);
  \draw[->] (0,0) -- (0,4.3);
  \foreach \t/\n in {0/1,1/2,2/4,3/8,4/16,5/32,6/64}
    \node[below, inner sep=1.5pt] at (\t,0) {\n};
  \foreach \v in {0,...,4}
    \node[left, inner sep=1.5pt] at (0,\v) {\v};
  \node[below] at (3,-0.45) {\iflangja{コア数 $N$}{cores $N$}};
  \node[rotate=90, above] at (-0.45,2) {\iflangja{相対性能}{relative performance}};
  \draw[thick, ln-muted, dashed] plot[domain=0:3, samples=30, smooth]
    (\x, {2^(2*\x/3)});
  \node[ln-muted, right, inner sep=1pt] at (3.15,3.7) {\iflangja{全並列}{fully parallel}};
  \draw[thick, ln-accent] plot[domain=0:6, samples=60, smooth]
    (\x, {2^(-\x/3)/(0.1+0.9/2^\x)});
  \node[ln-accent, anchor=north, inner sep=1pt] at (5.0,1.95) {$f_{\mathrm{par}}=0.9$};
  \draw[thick, ln-caution] plot[domain=0:6, samples=40, smooth]
    (\x, {2^(-\x/3)});
  \node[ln-caution, anchor=south west, inner sep=1pt] at (0.1,0.08) {\iflangja{単一スレッド}{single-threaded}};
\end{tikzpicture}
}
  \caption{Performance of $N$ cores sharing the power budget of one core,
    by \cref{eq:l22-perf}.}
  \label{fig:l22-power-perf}
\end{marginfigure}
Since the CPI does not change, a single-threaded program runs
$\sqrt[3]{N}$ times slower on a chip with $N$ cores than on a single core
with the same budget: more cores make a sequential program slower.  A
program whose parallel fraction is $f_{\mathrm{par}}$ gains from the
additional cores, as Amdahl's law (\cref{eq:l18-amdahl-cores}) describes,
but loses from the lower frequency:
\begin{equation}
  \text{performance}(N) \;=\; \frac{1}{\sqrt[3]{N}}\cdot
  \frac{1}{(1-f_{\mathrm{par}}) + \dfrac{f_{\mathrm{par}}}{N}} ,
  \label{eq:l22-perf}
\end{equation}
relative to the single core.  More generally, if a fraction $t_{p}$ of the
single-core execution time can use $p$ threads, the program needs
$\sqrt[3]{N}\,\sum_{p} t_{p}/\min(p,N)$ of that time on $N$ cores.
Beyond some number of cores the loss in frequency outweighs the gain in
parallelism (\cref{fig:l22-power-perf}).

\begin{example}
  \label{ex:l22-power}
  A chip has a \SI{120}{W} budget; a single core using all of it runs at
  \SI{4.2}{GHz}.  Eight cores receive \SI{15}{W} each and run at
  $\SI{4.2}{GHz}/\sqrt[3]{8}=\SI{2.1}{GHz}$, so a single-threaded program
  runs at half its speed on the single core.  A program with
  $f_{\mathrm{par}}=0.9$ performs $0.5/(0.1+0.9/8)\approx 2.35$ times as well
  as on the single core.  With 27 cores at \SI{1.4}{GHz} it reaches
  $(1/3)/(0.1+0.9/27)=2.5$, but with 64 cores at \SI{1.05}{GHz} only
  $0.25/(0.1+0.9/64)\approx 2.19$.
\end{example}

\section{Turbo mode}
\label{sec:l22-turbo}

When a program has no parallelism, the power budget of the idle cores can
be given to the core that runs it.\source{Lesson 22, video 15; H\&P §1.5.}
In \term{turbo mode}[ターボモード] the chip keeps only one core (or a few)
active, turns the others off, and raises the voltage and the clock frequency
of the active core above their normal values.  With the entire budget of $N$
cores, the frequency could rise by a factor of $\sqrt[3]{N}$---1.59 for four
cores---which would recover exactly the frequency lost by dividing the
budget (\cref{eq:l22-freq}).  \Cref{tab:l22-turbo} shows the two
processors of 2013 that the lecture uses as examples; their design power is
that of the whole package, which also contains the integrated graphics and
the uncore, and is therefore not a budget for the cores alone.

\begin{table}[htbp]
  \centering\small
  \caption{Normal and turbo clock frequencies of two quad-core Intel Core
    processors.  The power factor is the cube of the ratio of the turbo
    frequency to the normal one.}
  \label{tab:l22-turbo}
  \begin{tabular}{@{}lccccc@{}}
    \toprule
    Processor & Design power & Normal & Turbo & Frequency & Power \\
    \midrule
    i7-4702MQ (laptop) & \SI{37}{W} & \SI{2.2}{GHz} & \SI{3.2}{GHz} & $1.45\times$ & $3.1\times$ \\
    i7-4771 (desktop)  & \SI{84}{W} & \SI{3.5}{GHz} & \SI{3.9}{GHz} & $1.11\times$ & $1.38\times$ \\
    \bottomrule
  \end{tabular}
\end{table}

Neither processor uses the whole budget of four cores for one core.  What
limits a chip is its temperature, not only its total power: in turbo mode
the power is concentrated in the area of one core, only part of its heat
spreads into the idle cores around it, and the active core reaches its
maximum temperature before it consumes four times its normal power.  The
laptop processor normally runs its cores at low power, so one core can
still raise its power about threefold.  The desktop processor, designed for
better cooling, runs its cores at high power already in normal operation,
close to what the area of one core can dissipate, and gains little from
turbo mode.

\section{SMT, cores and chips}
\label{sec:l22-os}

The last challenge concerns the operating system.\source{Lesson 22, videos
16--17.}  A many-core system uses all the kinds of parallelism of Lesson 18
at once: a board carries several chips, each chip several cores, and each
core runs several threads with SMT.  The operating system sees every
hardware thread as a \term{logical processor}[論理プロセッサ] on which it
can run one software thread; two chips with eight cores each and two SMT
threads per core, for example, appear as 32 logical processors.

The logical processors are not equivalent, however.  Two threads on the same
core share its pipeline and caches, and each runs much slower than on a core
of its own (\cref{sec:l18-mt-performance}).  Threads on the same chip share
its last-level cache, its power budget and its memory interface.  Threads on
different chips share none of these, but communicate over the slower links
between the chips.  An operating system that treats all logical processors
as identical may start three threads on logical processors 0, 1 and 2,
which places two of the threads on one core and leaves the other chip idle
(\cref{fig:l22-placement}a).  A topology-aware operating system spreads the
threads first over the chips, then over the cores of each chip, and uses the
second SMT thread of a core only when every core is busy
(\cref{fig:l22-placement}b).

\begin{figure}[htbp]
  \centering
  % Two chips with two cores each and two SMT threads per core give eight
% logical processors (0-7).  Three threads A, B, C placed (a) on logical
% processors 0-2 without regard to the topology, (b) on separate chips and
% cores.
\begin{tikzpicture}[x=1cm,y=1cm, font=\scriptsize\sffamily,
  chip/.style={draw, fill=ln-box, rounded corners=2pt},
  core/.style={draw=ln-muted, fill=white},
  lp/.style={draw, minimum size=0.42cm, inner sep=0pt, font=\footnotesize\sffamily\bfseries},
  busy/.style={fill=ln-accent!25},
  num/.style={font=\tiny\sffamily, text=ln-muted, inner sep=1.5pt},
  paneltitle/.style={font=\footnotesize\sffamily, anchor=north}]
  \foreach \panel/\px/\captext in {0/0/{\iflangja{(a) トポロジを考慮しない配置}{(a) topology-unaware}},
                           1/5.95/{\iflangja{(b) トポロジを考慮した配置}{(b) topology-aware}}} {
    \begin{scope}[xshift=\px cm]
      \foreach \chipi in {0,1} {
        \pgfmathsetmacro{\chx}{\chipi*2.8}
        \draw[chip] (\chx,0) rectangle (\chx+2.55,1.7);
        \node[anchor=north west, inner sep=2pt] at (\chx,1.7) {\iflangja{チップ}{chip}~\chipi};
        \foreach \corei in {0,1} {
          \pgfmathsetmacro{\cox}{\chx+0.15+\corei*1.2}
          \pgfmathtruncatemacro{\coid}{2*\chipi+\corei}
          \draw[core] (\cox,0.1) rectangle (\cox+1.05,1.33);
          \node[anchor=north west, inner sep=1.5pt, text=ln-muted] at (\cox,1.33) {\iflangja{コア}{core}~\coid};
          \foreach \smt in {0,1} {
            \pgfmathtruncatemacro{\lpid}{4*\chipi+2*\corei+\smt}
            \node[lp] (lp-\panel-\lpid) at (\cox+0.3+\smt*0.45,0.7) {};
            \node[num, below] at (lp-\panel-\lpid.south) {\lpid};
          }
        }
      }
      \node[paneltitle] at (2.675,-0.05) {\captext};
    \end{scope}
  }
  % thread placements
  \foreach \panel/\lpid/\thr in {0/0/A, 0/1/B, 0/2/C, 1/0/A, 1/4/B, 1/2/C} {
    \node[lp, busy] at (lp-\panel-\lpid) {\thr};
  }
\end{tikzpicture}

  \caption{Three threads A, B and C on two chips with two cores and two SMT
    threads per core; small numbers are logical processors.  (a) Threads A
    and B share core~0.  (b) Every thread has a core of its own.}
  \label{fig:l22-placement}
\end{figure}

\begin{note}
  Spreading is best for independent threads; threads that share much data
  communicate faster on one chip, where the data stays in one LLC.  Either
  choice requires that the operating system know the topology.
\end{note}

\section{Challenges and solutions}
\label{sec:l22-summary}

\Cref{tab:l22-challenges} collects the challenges of this lesson and their
solutions.\source{Lesson 22, videos 18--19.}  Each arises because a
mechanism that serves a few cores does not grow with the number of cores.

\begin{table}[htbp]
  \centering\small
  \caption{Many-core challenges and their solutions.}
  \label{tab:l22-challenges}
  \begin{tabular}{@{}l>{\raggedright\arraybackslash}p{0.36\linewidth}>{\raggedright\arraybackslash}p{0.33\linewidth}@{}}
    \toprule
    Challenge & Cause & Solution \\
    \midrule
    Coherence traffic & more writes to shared data; a bus carries one
      transfer at a time & network on chip (mesh), directory-based coherence \\
    Off-chip traffic & more cache misses; pins grow more slowly than cores
      & distributed LLC \\
    Directory size & an entry for every memory block & on-chip partial
      directory \\
    Power budget & one budget shared by more cores & turbo mode when few
      cores are active \\
    Thread placement & SMT threads, cores and chips differ &
      topology-aware operating system \\
    \bottomrule
  \end{tabular}
\end{table}

\begin{keypoint}[Lesson summary]
  A many-core processor cannot simply scale up the designs of Lessons
  18--21.  Coherence traffic outgrows a bus, so the cores are connected by a
  network on chip, usually a mesh of tiles whose capacity, unlike that of a
  bus, grows with the number of cores---with $\sqrt{N}$ for traffic across
  the chip and with $N$ for traffic between neighbours---and use
  directory-based coherence.  A distributed LLC gives each tile a slice,
  selected by cache index or, for locality, by page number; the coherence
  directory becomes an on-chip partial directory in the tile of each block's
  LLC slice.  A shared power budget lowers the frequency of every core by a
  factor of $\sqrt[3]{N}$ (\cref{eq:l22-freq}), which turbo mode partly
  recovers for sequential programs, and the operating system must know which
  logical processors share a core or a chip.
\end{keypoint}

\end{document}
